%-------------------------------------------------------------------------------
%	SECTION TITLE
%-------------------------------------------------------------------------------
\cvsection{Projeler}

{\small\itshape\color{graytext}Farklı roller kapsamında yürütülen çok sayıda projeden seçilmiş öne çıkan girişimler.\par}
\vspace{-1mm}

%-------------------------------------------------------------------------------
%	CONTENT
%-------------------------------------------------------------------------------
\begin{cventries}

%---------------------------------------------------------
  \cventry
    {PROJE LİDERİ} % Role
    {Enerji İzleme, Tahminleme ve Anomali Tespiti Platformu} % Event
    {Manisa, Türkiye} % Location
    {2025 -- Devam Ediyor} % Date(s)
    {
      \begin{cvitems} % Description(s)
        \item {\textbf{200+} elektrik, su ve doğal gaz sayacını uzaktan okumalı sayaçlara çevirdi, altyapı veritabanını kurdu ve makine bazlı geçmiş enerji tüketimini, birim ürün başına enerji tüketimini ve karbon ayak izini görselleştiren bir enerji dashboard'u geliştirdi --- \textbf{SAP} ile entegre ederek tesisi \textbf{ISO 50001} sertifikasyonuna hazırladı.}
        \item {Kompresörden alınan basınçlı hava verisi, \textbf{MES}'ten gelen anlık ürün bilgisi ve \textbf{SAP}'den makine bazlı üretim teyit verilerini birleştiren bir hava kaçağı anomali tespit algoritması geliştirdi; erken müdahale ile kompresör enerji tüketimini azalttı.}
        \item {Planlanan üretim miktarlarına dayalı ML tabanlı enerji tahminlemesi ekledi (\textbf{XGBoost}, LightGBM, GBM regresyonu, R\textsuperscript{2} değeri \textbf{0.944}'e kadar); bu çalışma 13. Uluslararası Avrupa Disiplinlerarası Bilimsel Araştırmalar Konferansı'nda hakemli araştırma olarak yayınlandı (Tiran, Arnavutluk, Mayıs 2026).}
      \end{cvitems}
    }

%---------------------------------------------------------
  \cventry
    {PROJE LİDERİ } % Role
{Maestro} % Event
    {Manisa, Türkiye} % Location
    {Ara 2023 – Devam Ediyor} % Date(s)
    {
      \begin{cvitems} % Description(s)
        \item {\textbf{MES}, \textbf{SAP}, \textbf{PLC}, \textbf{IIoT} ve 3. parti verilerini fabrika genelinde birleştiren bir \textbf{Unified Namespace (UNS)} kurdu; \textbf{20 profesyonelden} oluşan bir ekibe liderlik etti. \textbf{3 fabrikada} yeni OT ağ altyapısını (NAT cihazları, OT switch'ler, firewall) kurdu ve \textbf{RBAC}/\textbf{LDAP} ile güvenli bir veri mimarisi tasarladı; izlenebilirliği \textbf{\%40} artırdı.}
        \item {\textbf{Power BI} raporları ve gerçek zamanlı makine verisi, birim ürün bazlı test verileri ile gerçek zamanlı \textbf{SPC} tabanlı anomali uyarıları gösteren özel dashboard'lar ile uçtan uca izlenebilirlik sağladı --- fabrikanın dijital ikizini güçlendirdi ve Digital Product Passport'unu besledi; \textbf{Kafka}, \textbf{Node.js} ve \textbf{React} kullanarak raporlama süresini \textbf{\%60} hızlandırdı, yıllık \textbf{\$156K} tasarruf sağladı.}
        \item {\textbf{YOLO} tabanlı iki bilgisayarlı görü güvenlik modeli entegre etti: robot çalışma alanına insan girdiğinde otomatik olarak durduran \textbf{AI-ODS} ve kör noktalarda forklift-yaya ile forklift-forklift çarpışmalarını hologram tarzı uyarılarla önleyen \textbf{AI-FDS} --- sıfır iş kazasına katkı sağladı.}
        \item {Titreşim ve sıcaklık sensörleri kullanarak kestirimci bakım uygulaması ekledi; \textbf{SAP} bağlantılı ekipman ömür takibi yaptı.}
      \end{cvitems}
    }

%---------------------------------------------------------
  \cventry
    {PROJE LİDERİ} % Role
    {Akıllı Sulama Hidrantı Sistemi (TÜBİTAK TEYDEB 1501)} % Event
    {Manisa, Türkiye} % Location
    {Oca 2022 – Kas 2022} % Date(s)
    {
      \begin{cvitems} % Description(s)
        \item {\textbf{TÜBİTAK TEYDEB 1501}'in \textbf{\textasciitilde3,5M₺} bütçesinden sorumlu proje yürütücüsü olarak, akıllı sulama hidrantı sisteminin ön yükleme ve ultrasonik sayaç modülleri için gömülü yazılım ile masaüstü/mobil uygulamalar geliştiren \textbf{9+} kişilik mühendislik ekibine liderlik etti; dışarıdan tedarik edilen çözüme kıyasla birim maliyeti \textbf{\textasciitilde\%68} azalttı ve satış sonrası teknik desteği tamamen şirket içine taşıdı.}
      \end{cvitems}
    }

%---------------------------------------------------------
\end{cventries}
